% Core idea and method description for the CoT-UQ BD paper.
% This file is read by the /paper skill as grounding for the Methodology section.

\section*{Core Idea (Internal Notes)}

\subsection*{Motivation}

Existing black-box uncertainty quantification (UQ) for LLMs measures confidence by
checking whether \emph{final answers} agree across multiple sampled responses.
This ignores the intermediate reasoning structure that Chain-of-Thought (CoT) models
explicitly produce. Two chains that reach the \emph{same final answer} via
\emph{different reasoning paths} represent genuine epistemic uncertainty that
answer-level methods discard entirely. Our methods exploit this signal.

\subsection*{Mathematical Foundation: Band Depth}

Band depth is a functional data analysis concept that ranks a set of curves from most
central (highest depth) to most outlying (lowest depth). The curve with the highest
band depth is the functional median. We extend this concept from ordered curves to
discrete reasoning paths on a graph: a CoT chain has high band depth if it lies
``inside'' the convex hull formed by other chains.

\subsection*{Method 1: pBD --- Path Band Depth}

Each CoT chain is treated as an \emph{ordered sequence of steps}. Chains of
different lengths are aligned via Dynamic Time Warping (DTW). The centrality of
chain $p_i$ is measured by the mean cosine similarity between each of its
step embeddings and the ensemble centroid at the corresponding aligned position
across all other chains:
\[
s_i = \frac{1}{n_i} \sum_{k=1}^{n_i}
\cos\!\left(e^{p_i}(k),\; \frac{1}{M-1}\sum_{r \neq i}
\overline{e^{p_r}(\sigma_r(k))}\right)
\]
where $\sigma_r$ is the DTW alignment from chain $p_i$ to reference chain $p_r$.
Confidence = $\bar{s}$ (mean depth over all $M$ chains).

\textbf{Limitation:} Order-sensitive --- chains $A \to B \to C \to D$ and
$A \to C \to B \to D$ are treated as different even if they express the same
reasoning. Also requires DTW alignment, which is $O(n^2)$ per pair.

\subsection*{Graph Construction (shared by vBD and gBD)}

From the $M$ sampled CoT chains, embed every reasoning step using a sentence
encoder. Build a reasoning graph $G = (V, E)$ as follows:
\begin{itemize}
  \item Steps whose cosine similarity exceeds a high threshold $\theta_c$ are merged
        into the same node (semantic clustering).
  \item Consecutive steps within the same chain are connected by a directed edge.
  \item Steps from different chains whose cosine similarity exceeds a lower threshold
        $\theta_e$ are connected by an undirected edge.
\end{itemize}
The result is a reasoning graph over the full ensemble where nodes represent
semantic clusters of reasoning states and edges represent either sequential
or semantic proximity relationships.

\subsection*{Method 2: vBD --- Vertex Band Depth}

Band depth is computed at the \emph{vertex} level on the reasoning graph.
For vertex $v$, the vertex band depth is the probability that $v$ lies inside
the geodesic convex hull of two randomly sampled vertices:
\[
v\text{BD}(v) = \mathbb{E}_{\mathcal{V}_2}\bigl[\chi(v \in H[\mathcal{V}_2])\bigr]
\]
The depth of chain $P_i$ is the mean vertex depth over all nodes it visits:
\[
c\text{BD}(P_i) = \frac{1}{|P_i|} \sum_{v \in P_i} v\text{BD}(v)
\]
\textbf{Advantage:} Order-invariant (only the \emph{set} of visited nodes
matters, not their sequence).
\textbf{Limitation:} Coarse --- treats each vertex independently and loses
path-level structural information.

\subsection*{Method 3: gBD --- Graph Band Depth (best)}

Band depth is restored to the \emph{path} level, but the reference distribution
is \textbf{random walks on the reasoning graph} rather than the $M$ sampled chains.
For each real CoT chain $P_i$, we generate $T$ pairs of random walks
$(P_a^{(t)}, P_b^{(t)})$ on $G$ and measure the fraction of $P_i$'s vertices
that fall inside their geodesic convex hull (soft containment):
\[
\text{gBD}(P_i) = \frac{1}{T} \sum_{t=1}^{T}
\frac{|\{v \in P_i : v \in \text{hull}(P_a^{(t)}, P_b^{(t)})\}|}{|P_i|}
\]
Final confidence score:
\[
\text{conf} = \alpha \cdot \overline{\text{gBD}} + (1-\alpha) \cdot \text{majority\_frac},
\quad \alpha = 0.10
\]

\textbf{gBD\_v1} (anchored): random walks are constrained to start at the
input node and end at the output node of the CoT graph.

\textbf{gBD\_v2} (free, best): random walks may start and end anywhere in $G$,
freely exploring the full topology of the reasoning graph.

\textbf{Key advantage:} By sampling reference paths not present in the ensemble,
gBD characterises the \emph{full reasoning space} defined by the graph topology,
not just the $N$ specific chains the model happened to generate. This makes it
robust to small ensembles and gives a richer band depth estimate. It is also
order-invariant.

\subsection*{Comparison Summary}

\begin{tabular}{lccc}
\hline
& \textbf{pBD} & \textbf{vBD} & \textbf{gBD} \\
\hline
Order-sensitive & Yes & No & No \\
Graph required  & No  & Yes & Yes \\
Reference distribution & Ensemble chains & Ensemble vertices & Random walks on graph \\
Depth level     & Path (step-wise) & Vertex, avg over chain & Path (soft hull) \\
Extra LLM calls & No  & No  & No \\
\hline
\end{tabular}

All three methods are \textbf{model-free}: they require only a small sentence encoder
(all-MiniLM-L6-v2) and no additional LLM inference calls beyond the initial ensemble.
